\subsubsection{Crear y consumir notificaciones automáticas de eventos del sistema}

\textbf{Descripción:} Se generan notificaciones para: registro de usuario, aprobación, rechazo, finalización de análisis, aprobación de análisis y marcado para repetir. El usuario las lista, marca como leídas, cuenta no leídas y puede eliminarlas.

\vspace{0.5em}
\noindent\textbf{PreCondiciones:}
\begin{itemize}[nosep]
    \item Usuario autenticado.
    \item Evento disparado existente en el servicio.
\end{itemize}

\vspace{0.5em}
\noindent\textbf{PostCondiciones:}
\begin{itemize}[nosep]
    \item Notificación persistida y visible al destinatario; estado cambia a leída si se marca.
\end{itemize}

\vspace{0.5em}
\noindent\textbf{Actores:} Analista, Administrador.

\vspace{0.5em}
\noindent\textbf{Flujo de eventos principal:}
\begin{enumerate}[nosep]
    \item Ocurre evento (ej. finalizar análisis).
    \item Servicio identifica destinatarios (p.ej., administradores o creador).
    \item Crea notificación (tipo, mensaje, referencia, \texttt{leída=false}).
    \item Usuario solicita listado paginado o no leídas.
    \item Selecciona una; se marca leída.
    \item Opcional: marca todas como leídas o elimina una.
\end{enumerate}

\vspace{0.5em}
\noindent\textbf{Flujos alternativos:}

\vspace{0.3em}
\noindent\textit{A1. Eliminación de notificación inexistente}
\begin{itemize}[nosep]
    \item Lanza error y no cambia estado.
\end{itemize}

\vspace{0.3em}
\noindent\textit{A2. Error al crear (repositorio)}
\begin{itemize}[nosep]
    \item Se captura y continúa sin bloquear el flujo principal del evento.
\end{itemize}

\vspace{0.5em}
\noindent\textbf{Requerimientos especiales:}
\begin{itemize}[nosep]
    \item Tipos de notificaciones: \texttt{USUARIO\_REGISTRO}, \texttt{USUARIO\_APROBADO}, \texttt{USUARIO\_RECHAZADO}, \texttt{ANALISIS\_FINALIZADO}, \texttt{ANALISIS\_APROBADO}, \texttt{ANALISIS\_REPETIR}.
    \item Operaciones: crear, listar, obtener no leídas, marcar leída, marcar todas, eliminar, contar no leídas.
\end{itemize}
